%---------- Inleiding ---------------------------------------------------------

% TODO: Is dit voorstel gebaseerd op een paper van Research Methods die je
% vorig jaar hebt ingediend? Heb je daarbij eventueel samengewerkt met een
% andere student?
% Zo ja, haal dan de tekst hieronder uit commentaar en pas aan.

%\paragraph{Opmerking}

% Dit voorstel is gebaseerd op het onderzoeksvoorstel dat werd geschreven in het
% kader van het vak Research Methods dat ik (vorig/dit) academiejaar heb
% uitgewerkt (met medesturent VOORNAAM NAAM als mede-auteur).
% 

\section{Inleiding}%
\label{sec:inleiding}
Deze bachelorproef onderzoekt methoden om AI-gegenereerde of AI-gemanipuleerde profielfoto's te detecteren met behulp van deep learning. De casus is gericht op organisaties die grote aantallen gebruikersfoto's verwerken (bv. sociale media, datingsites) en behoefte hebben aan betrouwbare automatische detectie om misleiding en identiteitsfraude te beperken.

Doelgroep: IT-onderzoekers en engineers die werken aan beeldforensische detectiesystemen, en operationele teams die zulke detectors willen integreren in moderatie- of verificatiepipelines.

Probleemstelling en centrale onderzoeksvraag:
\begin{quote}
\textit{Hoe kan een automatische detectiepijplijn onderscheid maken tussen authentieke gezichten en AI-gegenereerde/gemanipuleerde profielfoto's, zodanig dat de oplossing robuust generaliseert naar ongeziene datasets en uitlegbaar blijft voor audit en foutanalyse?}
\end{quote}

Onderzoeksdoelstelling: het ontwerpen, implementeren en evalueren van een proof-of-concept deepfake-detector voor stilstaande beelden, inclusief cross-dataset evaluatie, robuustheidstests (compressie) en uitlegbaarheid (Grad-CAM++). Het eindresultaat omvat een technisch rapport, geëxperimenteerde modellen en reproduceerbare evaluatiescripts.

%---------- Stand van zaken ---------------------------------------------------

\section{Literatuurstudie}%
\label{sec:literatuurstudie}

De literatuur richt zich op drie relevante thema's: (1) evolutie van generatieve modellen (GAN → diffusion) en de implicaties voor forensische cues; (2) detectiebenaderingen: CNN-gebaseerde lokale detectors, Vision Transformers voor globale consistentie, en hybride modellen; (3) evaluatiepraktijken: cross-dataset validatie en testen van robuustheid. Recente studies tonen dat klassieke forensische cues minder betrouwbaar worden naarmate generators realistischer worden, waardoor data-gedreven, generaliseerbare modellen en strenge cross-dataset evaluatie cruciaal zijn (zie hoofdstuk \ref{ch:stand-van-zaken}).

Dit voorstel bouwt expliciet voort op de bevindingen uit de literatuur: combinaties van architecturen en uitgebreide validatiestrategieën verminderen de generalisatiekloof en verhogen praktisch nut.

% Voor literatuurverwijzingen zijn er twee belangrijke commando's:
% \autocite{KEY} => (Auteur, jaartal) Gebruik dit als de naam van de auteur
%   geen onderdeel is van de zin.
% \textcite{KEY} => Auteur (jaartal)  Gebruik dit als de auteursnaam wel een
%   functie heeft in de zin (bv. ``Uit onderzoek door Doll & Hill (1954) bleek
%   ...'')


%---------- Methodologie ------------------------------------------------------
\section{Methodologie}%
\label{sec:methodologie}

De methodologie volgt een experimenteel, iteratief traject in vier fasen, consistent met hoofdstuk \ref{ch:methodologie} van de scriptie:
\begin{enumerate}
    \item \textbf{Data en Preprocessing:} Samenstellen van trainingsset (DeepDetect-2025) en onafhankelijke testset (140k), dataset-analyse, class-balancing (WeightedRandomSampler) en augmentatie pipelines.
    \item \textbf{Modelontwikkeling:} Implementatie van een model-factory die EfficientNet-B3 (CNN), ViT-Base/16 (transformer) en ConvNeXt-Small (hybride) laadt, vervangt laatste laag en traint met AdamW + cosine-annealing.
    \item \textbf{Evaluatie en Robuustheid:} Standaard metrieken (Accuracy, Precision, Recall, F1, ROC-AUC, PR-AUC), cross-dataset evaluatie op 140k, JPEG-compressie experiments en threshold-optimalisatie.
    \item \textbf{Uitlegbaarheid en Foutanalyse:} Grad-CAM++ heatmaps, most-confident-error analyse en aanbevelingen voor productie-inzet.
\end{enumerate}

Tools en omgeving: implementatie in Python 3.10+, PyTorch, timm, scikit-learn. Trainingsomgeving: Google Colab (NVIDIA T4) met aandacht voor geheugen (dataloader fixes). Code en evaluatiescripts worden opgeslagen in de `code/` directory en zijn reproduceerbaar.

Deliverables per fase:
\begin{itemize}
  \item Fase 1: dataset-samenvattingen, preprocessing scripts, `bachproef/code/dataset\_analysis.py`.
  \item Fase 2: getrainde modellen (checkpoints), `bachproef/code/build\_model.py`, `bachproef/code/training\_loop.py`.
  \item Fase 3: evaluatierapporten, figuren (ROC/PR, confusion matrices), robuustheid plots.
  \item Fase 4: Grad-CAM figuren, foutanalyse en een set aanbevelingen voor deployment.
\end{itemize}

%---------- Verwachte resultaten ----------------------------------------------
\section{Verwacht resultaat, conclusie}%
\label{sec:verwachte_resultaten}

Verwachte resultaten:
\begin{itemize}
  \item Modellen die op de trainingsdata hoge F1-scores halen, met variatie tussen CNN/ViT/hybride afhankelijk van compressie en datasetkenmerken.
  \item Cross-dataset evaluatie die de mate van generalisatie kwantificeert: naar verwachting daalt performance, maar eval op 140k geeft inzicht in welke architectuur het meest robuust is.
  \item Grad-CAM heatmaps die inzicht geven in waar modellen op focussen bij correcte voorspellingen en bij most-confident errors.
  \item Aanbevelingen voor drempelinstelling en productiegebruik, inclusief suggesties voor kalibratie en monitoring.
\end{itemize}

Meerwaarde voor de doelgroep: een reproduceerbare evaluatiepipeline, concrete modelaanbevelingen en XAI-instrumenten waarmee operationele teams detecties kunnen auditen en verfijnen.

Motivatie: gezien de snelle vooruitgang in generatieve modellen is een empirische studie met strikte cross-dataset validatie essentieel om praktische bruikbaarheid te beoordelen.

% Optionele korte planning
\section{Tijdsplanning en mijlpalen}
\begin{itemize}
  \item Week 1--3: data-voorbereiding.
  \item Week 4--7: modelimplementatie en korte trainingsruns (baseline).
  \item Week 8--11: uitgebreide trainingsrondes, hyperparam-zoek, robuustheidstests.
  \item Week 12--14: Grad-CAM analyses, foutanalyse en finalisatie resultaten.
  \item Week 15--16: schrijven eindrapport, figuren en bijlagen.
\end{itemize}

Concrete deliverables: code repository met notebooks/scripts, model-checkpoints (indien opslag-ruimte), alle figuren voor de scriptie en de geschreven bachelorproef.

% Einde voorstel-inhoud

